%------------------------------------------------------------------------------
% \Referat
%------------------------------------------------------------------------------
\chapter*{\vspace*{-\baselineskip}\begin{center}РЕФЕРАТ\end{center}\vspace*{-\baselineskip*2}}
\thispagestyle{empty}

%------------------------------------------------------------------------------
% Text
%------------------------------------------------------------------------------
\hspace*{12.5 mm}IP-ЯДРО НЕЙРОННОЙ СЕТИ ПРЯМОГО РАСПРОСТРАНЕНИЯ ДЛЯ РАСПОЗНОВАНИЯ РУКОПИСНЫХ ЦИФР :
дипломный проект / Е. А. Кривальцевич. – Минск : БГУИР, 2025, – п.з. – 79 с.,
чертежей (плакатов) – 5 л. формата А1 и 2 л. формата А2.
\vspace{1\baselineskip}

Основной задачей данного дипломного проекта является аппаратная реализация 
нейронной сети, использующей слой двумерного разделимого преобразования.
В нем будет рассмотрен подход увеличение точности рапознования рукописных цифр 
при не значительном увеличении количества параметров нейронной сети.

Ключевые этапы разработки включают:

– анализ существующих решений;

– обучение нейронной сети;

– разработка структуры IP-блока нейронной сети;

– проектирование эталонной модели;

– аппаратная реализация и тестирование;

– технико-экономическое обоснование разработки IP-ядра нейронной сети.

Предложенное решение по увеличению точности распознавания направлено на 
оптимизацию аппаратных затрат нейронной сети.

В целом предложенное решение по совершенствованию полносвязных нейронных сетей 
позволит значительно повысить эффективность.

\vspace{1\baselineskip}
Ключевые слова: нейронная сеть, двумерное разделимое преобразование,
MNIST, ПЛИС.